% =====================================================================
%  תרגיל בית 1 --- אחרי הרצאה 1 (משוואות פרידמן, פתרונות, מרחקים)
%  קומפילציה: lualatex hw1.tex
% =====================================================================
\documentclass[11pt,a4paper]{article}
% =====================================================================
%  Shared preamble — lecture notes "Introduction to Cosmology" (9 h)
%  Compile:  lualatex cosmo_lectures.tex   (twice, for the TOC)
% =====================================================================
\usepackage[a4paper,top=2.2cm,bottom=2.2cm,left=2.2cm,right=2.2cm]{geometry}
\usepackage{fontspec}
\usepackage{amsmath,amssymb}
\usepackage[bidi=basic]{babel}
\babelprovide[import,main]{hebrew}
\babelprovide[import]{english}

% ---- fonts: Culmus (Overleaf/Linux) -> Windows David -> FreeSerif ----
\IfFontExistsTF{David CLM}{%
  \babelfont{rm}{David CLM}\babelfont{sf}{Simple CLM}\babelfont{tt}{Miriam Mono CLM}%
}{\IfFontExistsTF{David}{%
  \babelfont{rm}[ItalicFeatures={FakeSlant=0.15}]{David}\babelfont{sf}{Arial}\babelfont{tt}{Courier New}%
}{\IfFileExists{/mnt/c/Windows/Fonts/david.ttf}{%
  \babelfont{rm}[Path=/mnt/c/Windows/Fonts/,Extension=.ttf,BoldFont=davidbd,ItalicFont=david,
    ItalicFeatures={FakeSlant=0.15}]{david}%
  \babelfont{sf}[Path=/mnt/c/Windows/Fonts/,Extension=.ttf,BoldFont=arialbd]{arial}%
  \babelfont{tt}[Path=/mnt/c/Windows/Fonts/,Extension=.ttf]{cour}%
}{%
  \babelfont{rm}{FreeSerif}\babelfont{sf}{FreeSans}\babelfont{tt}{FreeMono}%
}}}
\babelfont[english]{rm}{Latin Modern Roman}
\babelfont[english]{tt}{Latin Modern Mono}

\usepackage{xcolor}
\usepackage[framemethod=default]{mdframed}
\usepackage{enumitem}
\usepackage{booktabs}
\usepackage{graphicx}
\usepackage{fancyhdr}
\usepackage{needspace}
\usepackage{listings}
\usepackage[hidelinks]{hyperref}
\lstset{basicstyle=\ttfamily\footnotesize,keywordstyle=\color{qblue}\bfseries,
  commentstyle=\color{ggreen}\itshape,frame=single,rulecolor=\color{black!30},
  backgroundcolor=\color{black!2},breaklines=true,showstringspaces=false,
  language=Python,columns=fullflexible}

\renewcommand{\labelitemi}{\ensuremath{\bullet}}
\renewcommand{\labelitemii}{\ensuremath{\circ}}
\setlist[itemize]{leftmargin=1.4em,itemsep=1pt,topsep=2pt}
\setlist[enumerate]{leftmargin=1.6em,itemsep=1pt,topsep=2pt}

% ---- colours (as in the original course) ----
\definecolor{qblue}{RGB}{20,60,120}
\definecolor{ggreen}{RGB}{20,100,70}
\definecolor{kpurple}{RGB}{95,45,110}
\definecolor{worange}{RGB}{170,90,10}

\newmdenv[linecolor=kpurple,linewidth=1pt,backgroundcolor=kpurple!4,
  innertopmargin=6pt,innerbottommargin=6pt,innerleftmargin=8pt,innerrightmargin=8pt,
  skipabove=7pt,skipbelow=7pt]{keybox}
\newmdenv[linecolor=worange,linewidth=1pt,backgroundcolor=worange!5,
  innertopmargin=6pt,innerbottommargin=6pt,innerleftmargin=8pt,innerrightmargin=8pt,
  skipabove=7pt,skipbelow=7pt]{warnbox}
\newmdenv[linecolor=ggreen,linewidth=1pt,backgroundcolor=ggreen!5,
  innertopmargin=6pt,innerbottommargin=6pt,innerleftmargin=8pt,innerrightmargin=8pt,
  skipabove=7pt,skipbelow=7pt]{goalbox}
\newmdenv[linecolor=qblue,linewidth=1.1pt,backgroundcolor=qblue!5,
  innertopmargin=6pt,innerbottommargin=6pt,innerleftmargin=8pt,innerrightmargin=8pt,
  skipabove=8pt,skipbelow=6pt]{exbox}
\newmdenv[linecolor=black!35,linewidth=0.7pt,backgroundcolor=black!3,
  innertopmargin=6pt,innerbottommargin=6pt,innerleftmargin=8pt,innerrightmargin=8pt,
  skipabove=4pt,skipbelow=8pt]{solbox}

% ---- macros ----
\newcommand{\en}[1]{\foreignlanguage{english}{#1}}
\newcommand{\Msun}{M_{\odot}}
\newcommand{\Mpc}{\mathrm{Mpc}}
\newcommand{\kpc}{\mathrm{kpc}}
\newcommand{\Gyr}{\mathrm{Gyr}}
\newcommand{\kms}{\mathrm{km\,s^{-1}}}
\newcommand{\kmsMpc}{\mathrm{km\,s^{-1}\,Mpc^{-1}}}
\newcommand{\LCDM}{\Lambda\mathrm{CDM}}

% lecture header: \lecture{number}{title}{time budget + one-line content}
\newcounter{lec}
\counterwithin*{subsection}{lec}
\renewcommand{\thesubsection}{\thelec.\arabic{subsection}}
\makeatletter\renewcommand{\fps@figure}{htbp}\makeatother
\newcommand{\lecture}[3]{%
  \clearpage\refstepcounter{lec}%
  \noindent\colorbox{qblue}{\parbox{\dimexpr\textwidth-2\fboxsep}{%
   \color{white}\bfseries\Large הרצאה #1: #2\\[3pt]\normalsize\mdseries #3}}%
  \addcontentsline{toc}{section}{הרצאה #1: #2}\par\medskip}

% exercises: numbered per lecture (1.1, 1.2, ...)
\newcounter{ex}[lec]
\newenvironment{exercise}[1]{%
  \refstepcounter{ex}\begin{exbox}\textbf{\textcolor{qblue}{תרגיל \thelec.\theex\ --- #1.}}\ }%
  {\end{exbox}}
\newcommand{\soln}[1]{\par\smallskip\noindent\textbf{פתרון #1.}\ }

\pagestyle{fancy}\fancyhf{}
\fancyhead[R]{\small מבוא לקוסמולוגיה --- רשימות הרצאה}
\fancyhead[L]{\small\hbox{\textdir TLT\relax\thepage}}
\renewcommand{\headrulewidth}{0.4pt}
\setlength{\parindent}{0pt}\setlength{\parskip}{4pt}

% ---- shared definitions for the homework sheets (input after ../lecture_notes/preamble) ----
\newcounter{hwq}
\newenvironment{question}[2]{%
  \refstepcounter{hwq}\needspace{5\baselineskip}%
  \begin{exbox}\textbf{\textcolor{qblue}{שאלה \thehwq\ --- #1}}\hfill\textbf{(#2 נק')}\par\smallskip}%
  {\end{exbox}}
\newenvironment{parts}{\begin{enumerate}[label=(\alph*),itemsep=3pt]}{\end{enumerate}}
\newcommand{\hwtitle}[3]{%
  \noindent\colorbox{qblue}{\parbox{\dimexpr\textwidth-2\fboxsep}{%
   \color{white}\bfseries\Large תרגיל בית #1: #2\\[3pt]\normalsize\mdseries #3}}\par\medskip}
\newcommand{\hwrules}[1]{%
\begin{goalbox}\small\textbf{הנחיות.}
\begin{itemize}
\item ארבע שאלות, 25 נקודות כל אחת. השאלות בלתי תלויות זו בזו.
\item אלא אם נאמר אחרת: $\LCDM$ שטוח עם $h=0.674$, $\Omega_m=0.315$, $\Omega_\Lambda=1-\Omega_m$ (קרינה מוזנחת), ו-$c/H_0=4448\,\Mpc$, $1/H_0=14.51\,\Gyr$.
\item רוב הסעיפים אנליטיים: גזירה, ואז הצבת מספרים במחשבון. רק סעיפים המסומנים \textbf{[מחשב]} דורשים קוד; בהם צרפו את הקוד (\en{Python} מומלץ), סמנו צירים ויחידות, ושמרו את ה-\en{seed}.
\item כל טענה, נתון או ערך שלא נלמדו בהרצאה --- עם ציטוט (מחבר, שנה, כתב עת).
#1
\end{itemize}\end{goalbox}}
\pagestyle{fancy}\fancyhf{}
\fancyhead[L]{\small\hbox{\textdir TLT\relax\thepage}}

\fancyhead[R]{\small מבוא לקוסמולוגיה --- תרגיל בית 1}
\begin{document}
\hwtitle{1}{משוואות פרידמן ופתרונותיהן}{אחרי הרצאה 1 \quad$\cdot$\quad משוואת הנוזל, פתרונות לרכיב אחד ולשני רכיבים, עקמומיות, מרחקים}
\hwrules{}

% ---------------------------------------------------------------------
\begin{question}{משוואת הנוזל ומשוואת התאוצה}{25}
\begin{parts}
\item \textbf{(5)} נפח נע $V\propto a^3$ מתפשט אדיאבטית. מהחוק הראשון $d(\rho c^2V)=-p\,dV$ גזרו את משוואת הנוזל $\dot\rho+3H(\rho+p/c^2)=0$. הראו שעבור $p=w\rho c^2$ עם $w$ קבוע מתקבל $\rho\propto a^{-3(1+w)}$, ורשמו את $w$ ואת התלות ב-$a$ עבור חומר, קרינה וקבוע קוסמולוגי. הסבירו במשפט אחד למה קרינה מדוללת מהר יותר מחומר.
\item \textbf{(6)} גזרו בזמן את משוואת פרידמן הראשונה $H^2=\frac{8\pi G}{3}\rho-\frac{kc^2}{a^2}$, השתמשו במשוואת הנוזל, והראו שמתקבלת משוואת התאוצה
$\displaystyle\frac{\ddot a}{a}=-\frac{4\pi G}{3}\Big(\rho+\frac{3p}{c^2}\Big)$.
עבור איזה $w$ היקום מאיץ? למה לחץ \emph{חיובי} מאט את ההתפשטות, למרות שבגז רגיל לחץ "דוחף החוצה"?
\item \textbf{(6)} \emph{גזירה ניוטונית.} גלקסיה בשפת כדור אחיד בעל צפיפות $\rho$ נמצאת ברדיוס $R(t)=a(t)\,r$. כתבו את שימור האנרגיה $\tfrac12\dot R^2-GM/R=E$ והראו שהוא נותן את משוואת פרידמן הראשונה. מה הקשר בין $E$ ל-$k$? איזה איבר במשוואות פרידמן אין לו מקבילה ניוטונית?
\item \textbf{(8)} חשבו את $\rho_{c,0}=3H_0^2/8\pi G$ ביחידות $\mathrm{kg\,m^{-3}}$, בפרוטונים למטר מעוקב ובמסות שמש ל-$\Mpc^3$. עם $\Omega_bh^2=0.0224$, כמה בריונים יש במטר מעוקב ממוצע ביקום היום? השוו לצפיפות האוויר בחדר ($\sim10^{25}$ מולקולות ל-$\mathrm m^3$) ולצפיפות הגז הבין-כוכבי בשביל החלב ($\sim10^{6}\,\mathrm m^{-3}$).
\end{parts}
\end{question}

% ---------------------------------------------------------------------
\begin{question}{פתרונות שטוחים: רכיב אחד ושני רכיבים}{25}
\begin{parts}
\item \textbf{(6)} ליקום שטוח עם רכיב יחיד בעל $w\neq-1$ קבוע, הראו ש-$a\propto t^{2/[3(1+w)]}$, ומצאו את $H(t)$ ואת הגיל $t_0$ ביחידות $1/H_0$. הציבו לחומר ולקרינה. מה הפתרון עבור $w=-1$? למה בו אין "מפץ גדול"?
\item \textbf{(8)} ליקום שטוח עם חומר ו-$\Lambda$, הראו (בהצבה ישירה במשוואת פרידמן) ש-
\[a(t)=\Big(\frac{\Omega_m}{\Omega_\Lambda}\Big)^{1/3}\sinh^{2/3}\!\Big(\tfrac32\sqrt{\Omega_\Lambda}\,H_0t\Big)\]
הוא פתרון עם $a(0)=0$. בדקו את הגבולות $t\to0$ ו-$t\to\infty$: לאילו פתרונות מסעיף (א) הם מתכנסים?
\item \textbf{(6)} מצאו מהפתרון ביטוי סגור לגיל $t_0$ (השתמשו ב-$a(t_0)=1$), וחשבו אותו ב-$\Gyr$ עבור הפרמטרים של \en{Planck}. השוו לגיל של יקום \en{Einstein--de Sitter} עם אותו $H_0$ ולגיל הכוכבים העתיקים ביותר ($\approx13$ מיליארד שנה). למה זו הייתה בעיה בשנות ה-90?
\item \textbf{(5)} בעתיד הרחוק $H\to H_\infty$. מצאו את $H_\infty$ ב-$\kmsMpc$ ואת זמן ה-\en{e-folding} $1/H_\infty$ ב-$\Gyr$. בכמה יגדל המרחק בין שני צבירים רחוקים ב-100 מיליארד השנים הבאות (בקירוב של $H=H_\infty$)?
\end{parts}
\end{question}

% ---------------------------------------------------------------------
\begin{question}{עקמומיות וגורל היקום}{25}
עכשיו $\Omega_k=1-\Omega_m-\Omega_\Lambda$ ו-$H^2=H_0^2\big[\Omega_ma^{-3}+\Omega_ka^{-2}+\Omega_\Lambda\big]$.
\begin{parts}
\item \textbf{(9)} \emph{יקום סגור של חומר בלבד} ($\Omega_m>1$, $\Omega_\Lambda=0$). הראו שהפתרון הפרמטרי
\[a(\eta)=\frac{\Omega_m}{2(\Omega_m-1)}\,(1-\cos\eta),\qquad H_0t(\eta)=\frac{\Omega_m}{2(\Omega_m-1)^{3/2}}\,(\eta-\sin\eta)\]
מקיים את משוואת פרידמן (רמז: $dt=a\,d\eta/[H_0\sqrt{\Omega_m-1}]$). מצאו את $a_{\max}$ ואת זמן הקריסה $t_{\rm crunch}$.
\item \textbf{(6)} הציבו $\Omega_m=2$: מה $a_{\max}$? מה הגיל הנוכחי ומתי תהיה הקריסה (ב-$\Gyr$)? האם יקום כזה תואם את גיל הכוכבים?
\item \textbf{(5)} הראו שבלי $\Lambda$ היקום קורס אם ורק אם $\Omega_m>1$ (כלומר "גאומטריה = גורל"). רמז: בדקו מתי $E^2(a)=\Omega_ma^{-3}+(1-\Omega_m)a^{-2}$ מתאפס עבור $a>1$.
\item \textbf{(5)} הראו שעם $\Lambda>0$ הקשר הזה נשבר: ליקום הסגור $(\Omega_m,\Omega_\Lambda)=(0.3,0.8)$ מצאו אנליטית את המינימום של $E^2(a)$ והראו שהוא חיובי, כלומר היקום מתפשט לנצח. ולהפך: ביקום שטוח עם $\Lambda<0$, למשל $\Omega_m=1.1$, $\Omega_\Lambda=-0.1$, מצאו את $a_{\max}$.
\end{parts}
\end{question}

% ---------------------------------------------------------------------
\begin{question}{מרחקים, גדלים זוויתיים ומהירויות "על-אורית"}{25}
\begin{parts}
\item \textbf{(6)} \emph{יקום ריק (\en{Milne}).} עבור $\Omega_m=\Omega_\Lambda=0$ ($\Omega_k=1$) הראו ש-$a=H_0t$, ש-$\chi=\frac{c}{H_0}\ln(1+z)$, ש-$D_M=\frac{c}{H_0}\sinh\chi'$ (עם $\chi'=\ln(1+z)$) ולכן
$d_L=\frac{c}{H_0}\,z\big(1+\tfrac z2\big)$. בדקו את גבול $z\ll1$.
\item \textbf{(7)} גזרו את $d_A(z)$ בשלושה יקומים: \en{Milne} (מסעיף א), \en{Einstein--de Sitter}, ו-\en{de Sitter} שטוח ($\Lambda$ בלבד, $H=H_0$). שרטטו סקיצה של שלושתם. הראו ש-$d_A$ של \en{Milne} ושל \en{de Sitter} עולה מונוטונית (מה הגבול שלו ב-$z\to\infty$?), בעוד ל-\en{EdS} יש מקסימום (נזכיר: ב-$z=1.25$).
\item \textbf{(6)} מה הגודל הזוויתי (בשניות קשת) של גלקסיה בקוטר $10\,\kpc$ ב-$z=1$ וב-$z=3$ בשלושת המודלים? השוו ל"ציפייה האוקלידית" $\theta=10\,\kpc/(cz/H_0)$. באיזה מודל הגלקסיה ב-$z=3$ נראית גדולה יותר מאשר ב-$z=1$?
\item \textbf{(6)} ב-\en{EdS}: גלקסיה במרחק נע $D$ מתרחקת היום במהירות $v=H_0D$. הראו ש-$v=c$ ב-$z=3$. האם זו סתירה לתורת היחסות? האם אנו רואים גלקסיות כאלה? הראו שזמן המבט לאחור הוא $t_{\rm lb}=\frac{2}{3H_0}\big[1-(1+z)^{-3/2}\big]$, ושהיחס $D/(c\,t_{\rm lb})=3\frac{1-(1+z)^{-1/2}}{1-(1+z)^{-3/2}}$ גדול מ-1 לכל $z>0$. הציבו $z=1$ והסבירו את התוצאה.
\end{parts}
\end{question}

\end{document}
